\documentclass[11pt,a4paper]{article}
\usepackage[margin=2.5cm]{geometry}
\usepackage{amsmath,amssymb}
\usepackage{siunitx}
\usepackage{hyperref}
\usepackage[T1]{fontenc}
\usepackage{lmodern}

  \title{02\_Run\_Extraction: Straight-Line Coastdown Interval Detection (Legacy Technical Note)}
\author{Validation\_Aero\_Test}
\date{\today}

\begin{document}
\maketitle

\section{Overview}
This document describes the current straight-line interval detection used to extract coastdown runs.
The implementation lives in \texttt{aero\_validation/segmentation.py} and is exposed through the CLI script
  \texttt{Code/02\_Run\_Extraction/export\_drag\_runs.py}.

The algorithm:
\begin{itemize}
  \item proposes candidate run starts using speed peak detection,
  \item extends forward until a stopping condition (turning/braking/re-acceleration/stop),
  \item enforces quality gates (straightness, duration),
  \item writes an intervals CSV used by downstream steps.
\end{itemize}

\section{Inputs}
From the cleaned dataset CSV (typically \texttt{Outputs/01\_Data\_Preparation/Aero\_Validation\_Signals\_cleaned\_10ms.csv}):
\begin{itemize}
  \item \textbf{Time}: \texttt{t\_s} [s]
  \item \textbf{Speed}: \texttt{GPS1\_VEL\_MAG\_ms} [\si{\meter\per\second}]
  \item \textbf{Steering}: \texttt{ICU\_SteeringAngle\_deg} [deg]
  \item \textbf{Optional brakes}: \texttt{ICU\_Brakpresf\_bar}, \texttt{ICU\_Brakpresr\_bar} [bar]
\end{itemize}

\paragraph{Note}
This file is a legacy/technical note. For the maintained documentation, see
  \texttt{Docs/02\_Run\_Extraction/02\_Run\_Extraction.tex}.

\section{Preprocessing}
Let $t$ be the sample times and $v$ the speed signal.
\begin{itemize}
  \item Estimate the median sample period $\Delta t_\mathrm{med}$ from $\Delta t = \mathrm{diff}(t)$.
  \item Convert time windows (seconds) to sample windows by $n \approx \mathrm{round}(\Delta t_\mathrm{window}/\Delta t_\mathrm{med})$.
  \item If smoothing is enabled (default): smooth $v$ with a centered rolling mean using a window derived from \texttt{peak\_window\_s}.
  \item If brake signals exist: combine $b(t)=\max(b_f,b_r)$ and lightly smooth it with a centered rolling mean (fixed small window).
\end{itemize}

\section{Candidate peak detection}
Peak detection is implemented in \texttt{detect\_drag\_peaks(t, v, config)}.
It combines two mechanisms and returns sorted unique indices:
\begin{enumerate}
  \item \textbf{Local maxima}: indices where $v_s[i] \ge v_s[i-1]$ and $v_s[i] \ge v_s[i+1]$.
  \item \textbf{Derivative sign change}: indices where the speed derivative changes from clearly positive to non-positive.
\end{enumerate}

\paragraph{Prominence gate}
For each local maximum, compute a prominence proxy
\[
  p = v_s[i] - P_{60}(\,v_s\;\text{in a local window around }i\,),
\]
and require $v_s[i] \ge \texttt{min\_peak\_speed}$ and $p \ge \texttt{peak\_prominence}$.

\section{Forward extension and end conditions}
Given a start index $p$ (peak), the algorithm extends forward until the first end condition triggers.
Let $i$ be the current index and $j=i+\texttt{accel\_window\_s}$ (converted to samples).

End conditions (evaluated sequentially):
\begin{itemize}
  \item \textbf{Stop}: $v_s[i] \le \texttt{stop\_speed\_th}$.
  \item \textbf{Turning}: $|\theta[i]| > \texttt{steer\_thresh}$ sustained for \texttt{turn\_window\_s}.
  \item \textbf{Braking}: $b[i] > \texttt{brake\_thresh}$ or an instantaneous brake ``spike'' exceeds \texttt{brake\_spike\_thresh}.
  \item \textbf{Re-acceleration}: if $v_s[j] - v_s[i] > \texttt{reaccel\_eps}$.
\end{itemize}

The returned interval is half-open: $(p,e)$ where $e$ is the first index after the accepted end.

\section{Quality gates}
After computing $(p,e)$:
\begin{itemize}
  \item \textbf{Straightness}: the fraction of samples in $[p,e)$ where $|\theta| \le \texttt{steer\_thresh}$ must be at least \texttt{straight\_ratio}.
  \item \textbf{Duration}: the run duration $t[e-1]-t[p]$ must satisfy \texttt{min\_dur\_s} $\le$ duration $\le$ \texttt{max\_dur\_s}.
  \item \textbf{Non-overlap}: candidate peaks that would overlap a previously accepted interval are skipped.
\end{itemize}

\section{CLI usage and outputs}
\subsection*{Export intervals}
  \texttt{Code/02\_Run\_Extraction/export\_drag\_runs.py} loads the cleaned dataset, runs segmentation, writes an intervals CSV, and generates an overview plot.

Default output paths:
\begin{itemize}
  \item Intervals CSV: \texttt{Outputs/02\_Run\_Extraction/csv/drag\_intervals.csv}
  \item Overview plot: \texttt{Outputs/02\_Run\_Extraction/plots/velocity\_overview\_with\_runs.png}
\end{itemize}

Intervals CSV schema (key fields):
\begin{itemize}
  \item \texttt{run\_id}: 1..N
  \item \texttt{start\_idx}, \texttt{end\_idx}: index range into the input dataset (end exclusive)
  \item \texttt{t0\_s}, \texttt{t1\_s}, \texttt{duration\_s}
  \item \texttt{v\_start\_ms}, \texttt{v\_end\_ms}
  \item \texttt{source\_basename}, \texttt{row\_count}: metadata for mismatch detection
\end{itemize}

\subsection*{Optional debug log}
If \texttt{--debug-out} is provided, the debug variant \texttt{segment\_drag\_runs\_with\_debug} writes a per-peak decision log (accepted/rejected + reasons).

\section{Default thresholds (CLI defaults)}
The defaults below are the ones exposed by \texttt{export\_drag\_runs.py}:
\begin{center}
\begin{tabular}{ll}
            	    	\texttt{steer\_thresh} & 10 deg \\
            	    	\texttt{straight\_ratio} & 0.75 \\
            	    	\texttt{min\_peak\_speed} & 15.0 m/s \\
            	    	\texttt{peak\_window\_s} & 0.5 s \\
            	    	\texttt{peak\_prominence} & 0.05 m/s \\
            	    	\texttt{reaccel\_eps} & 0.2 m/s \\
            	    	\texttt{accel\_window\_s} & 0.5 s \\
            	    	\texttt{brake\_thresh} & 0.05 bar \\
            	    	\texttt{brake\_spike\_thresh} & 0.1 bar/s \\
            	    	\texttt{turn\_window\_s} & 0.3 s \\
            	    	\texttt{min\_dur\_s} & 1.5 s \\
            	    	\texttt{max\_dur\_s} & 20.0 s \\
            	    	\texttt{stop\_speed\_th} & 0.2 m/s \\
\end{tabular}
\end{center}

\section{Per-run outputs}
The script \texttt{Code/02\_Run\_Extraction/generate\_run\_outputs.py} consumes \texttt{drag\_intervals.csv} and writes:
\begin{itemize}
  \item \texttt{Outputs/02\_Run\_Extraction/csv/run\_\{XXX\}.csv}: sliced time series per interval
  \item \texttt{Outputs/02\_Run\_Extraction/plots/run\_\{XXX\}\_velocity\_steering\_brake.png}: QC plots per run
\end{itemize}

\end{document}
